\documentclass[11pt,a4paper]{article}
\usepackage[utf8]{inputenc}
\usepackage[T1]{fontenc}
\usepackage{geometry}
\geometry{margin=2.5cm}
\usepackage{amsmath,amssymb}
\usepackage{graphicx}
\usepackage{listings}
\usepackage{xcolor}
\usepackage{booktabs}
\usepackage{hyperref}
\hypersetup{colorlinks=true, linkcolor=blue!60!black, urlcolor=blue!60!black, citecolor=blue!60!black}

\lstset{
  language=Python,
  basicstyle=\ttfamily\footnotesize,
  keywordstyle=\color{blue!70!black}\bfseries,
  commentstyle=\color{green!45!black}\itshape,
  stringstyle=\color{red!60!black},
  frame=single,
  breaklines=true,
  captionpos=b,
  numbers=left,
  numberstyle=\tiny\color{gray},
  showstringspaces=false,
}

\title{A Non-Planar Slicer for Robot-Based FDM Printing\\of Continuous-Fiber Composites}
\author{}
\date{\today}

\begin{document}
\maketitle
\tableofcontents
\newpage

% ======================================================================
\section{Introduction}\label{sec:intro}
% ======================================================================

\subsection{Background and motivation}

Fused Deposition Modeling (FDM) is today the most widespread additive
manufacturing process for thermoplastics. In its standard form, a
heated nozzle extrudes a continuous filament that is deposited along a
toolpath, layer by layer, until the part is complete. Virtually every
commercial FDM machine and every mainstream slicing package
(Cura, PrusaSlicer, Slic3r, \dots) shares one structural assumption: the
print head moves on a three-axis Cartesian gantry, the nozzle axis
stays vertical, and every layer is a \emph{horizontal plane} parallel to
the build plate. The slicer reduces the solid model to a stack of 2D
regions; the kinematics of the printer reduces the tool orientation to
a single fixed pose.

This assumption is a software and kinematics limitation, not a physical
one. Nothing in the extrusion process requires the nozzle to be
vertical: material can be deposited on inclined or curved surfaces,
provided that (i) the tool can be oriented so that its axis is roughly
perpendicular to the target surface, and (ii) the deposited bead is
pressed against an existing substrate. A six-degree-of-freedom
industrial robot arm satisfies requirement (i) by construction: it can
place and orient an end-effector arbitrarily within its workspace,
with the repeatability needed for extrusion processes. Removing the
planar-layer constraint then unlocks concrete benefits:

\begin{itemize}
  \item \textbf{Elimination of the staircase effect.} Sloped and curved
        surfaces are no longer approximated by stacked planar slices;
        layers can follow the actual surface curvature, so the
        as-printed geometry matches the CAD geometry without the
        characteristic stair steps of planar printing.
  \item \textbf{Conformal reinforcement.} Continuous fibers can be laid
        along the actual load paths of a structural part, on curved
        surfaces, rather than being discretized into flat slices that
        cut through the mechanically optimal trajectories.
  \item \textbf{Printing on / against existing geometry.} A
        non-planar-capable robot can print support structures that
        conform to the underside of a part, or deposit material on
        pre-existing surfaces, which is the basis of hybrid
        manufacturing.
\end{itemize}

Realizing these benefits, however, requires a slicer that fundamentally
does more than a conventional one. A planar slicer only needs to output
2D coordinates and a feed rate, because the tool orientation is
constant. A robot-based non-planar slicer must output a complete
\emph{tool pose} --- position \emph{and} orientation --- at every point
of the trajectory, must guarantee that the deposited bead stays
supported by previously deposited material or by the build plate, and
must handle path topology (connections, gaps, crossings) that planar
slicers never encounter because they can always "travel" at a fixed
height.

\subsection{Project scope and objectives}

This project develops such a slicer: a Python software package, called
\texttt{slicer}, that transforms a structural finite-element model into
full six-DoF printing trajectories, ready to be simulated and executed
on a robotic cell. The application domain is the FDM printing of
\emph{continuous-fiber composite} parts: the target objects are
load-bearing structural components (the running example is a thin
C-beam) in which stiff carbon-fiber tows are embedded in a
thermoplastic matrix, and whose fiber layout is determined by a
structural optimization (compliance minimization) performed upstream.

Concretely, the slicer must:

\begin{enumerate}
  \item \textbf{Slice a solid model into non-planar layers} ---
        each layer is an isosurface of a scalar field defined on the
        volumetric mesh, so that layers naturally follow the curved
        geometry of thin-walled structural parts.
  \item \textbf{Generate the infill of each layer, aligned with the
        fiber trajectories} deposited inside that layer. This is the
        central mechanical requirement: the matrix material must support
        the fibers, and the fiber direction controls the local
        stiffness of the printed composite.
  \item \textbf{Generate contour shells} (offset perimeter passes) at a
        nozzle-width scale, on every layer, for both the part and its
        support structure.
  \item \textbf{Automatically generate a support structure (mold)}
        underneath the overhanging regions of the part, by computing a
        "shadow volume" and slicing it into printable layers with a
        grid infill, including the handling of infill crossings so
        that no bead is deposited twice.
  \item \textbf{Assemble every layer's curves into a continuous,
        physical printing path}: order the curves, connect them when the
        connection is short and legitimate, and insert perpendicular
        approach/leave moves when it is not.
  \item \textbf{Compute a full tool pose (rotation matrix) at every
        trajectory point} --- approach axis, travel axis and their
        cross product --- so that the exported trajectory can drive a
        six-axis robot.
  \item \textbf{Export the complete multi-layer trajectory} to a
        plain-text format that an external robot simulator (RoboDK) can
        ingest directly.
\end{enumerate}

\subsection{Inputs and outputs of the pipeline}

The slicer operates within a larger tool chain built around the
\texttt{oamc} library (Optimization of Anisotropic Material with
Continuous fibers). The data flow is:

\begin{itemize}
  \item \textbf{Input: an APDL (ANSYS) finite-element model}, parsed by
        \texttt{oamc.integrations.ansys.parser.APDLParser}. The model
        contains (a) a volumetric solid mesh of the part (linear
        tetrahedra, hexahedra, or their quadratic variants), (b) a
        named surface mesh \texttt{"MOLD"} describing the surface the
        part will be printed against, and (c) boundary conditions and
        material data (an isotropic matrix material and a transversely
        isotropic fiber material).
  \item \textbf{Input: fiber paths.} For each layer, a set of
        continuous fiber trajectories (polylines in 3D), either
        computed by the structural optimization of \texttt{oamc}
        (\texttt{model.compute\_fibers}) and saved as CSV files, or
        provided directly. Each file \texttt{fiber\_<layer>\_<i>.csv}
        contains one polyline; the slicer loads them per layer.
  \item \textbf{Output: printing trajectories.} One text file for the
        part, one for the mold, each containing, line by line, every
        trajectory point with its Cartesian coordinates, the nine
        components of its tool-orientation matrix, a segment-type label
        and the layer index.
\end{itemize}

The pipeline is deliberately kept as a pure geometry/trajjectory
generator: it does not simulate the robot, check collisions, or compute
extrusion rates. Those responsibilities belong to the downstream robot
simulator. This separation keeps the slicer generic (any six-DoF cell
can consume its output) and testable (the output format is trivially
inspectable).

\subsection{Software ecosystem}

Three open-source libraries carry the computational load:

\begin{itemize}
  \item \textbf{PyVista} \cite{pyvista}, a Pythonic interface to the
        Visualization Toolkit (VTK) \cite{vtk}. Every geometric
        operation of the slicer --- isosurface extraction, scalar-field
        clipping, contouring, slicing, feature-edge extraction, Delaunay
        triangulation, convex hulls, point merging --- is a VTK filter
        exposed by PyVista. This choice is central to the project: the
        whole slicing machinery reduces to compositions of
        well-tested, C++-implemented mesh filters, called from Python.
  \item \textbf{NumPy} for all array-level mathematics, and
        \textbf{SciPy} \cite{scipy} for spatial indexing
        (\texttt{cKDTree}), convex hulls (\texttt{ConvexHull}), Delaunay
        triangulation (\texttt{Delaunay}) and rotations
        (\texttt{spatial.transform.Rotation}).
  \item \textbf{RoboDK} \cite{robodk} as the downstream consumer of the
        exported trajectories, for offline programming and simulation of
        the robot cell.
\end{itemize}

\subsection{Report organization}

The report is organized as follows.
Chapter~\ref{sec:literature} reviews the state of the art of non-planar
slicing and multi-axis printing. Chapter~\ref{sec:methodology} presents
the methodology: the complete algorithm, stage by stage, from model
orientation to trajectory export, with the underlying mathematics.
Chapter~\ref{sec:implementation} details the Python implementation:
the class architecture, the role of each module and method, the
performance engineering, and the connection to RoboDK.
Chapter~\ref{sec:results} presents results (to be completed).

% ======================================================================
\section{Literature review}\label{sec:literature}
% ======================================================================

% ----------------------------------------------------------------------
% NOTE: this section is intentionally left as a placeholder.
% It will be written separately by the author.
% ----------------------------------------------------------------------

\begin{quote}
\itshape
\textbf{[To be completed.]} This chapter will review existing non-planar
slicing approaches: planar slicing and its limitations (staircase
effect), curved-layer slicing for FDM, conformal printing on
freeform surfaces, multi-axis / robotic FDM and full tool-pose
trajectory generation, and support-structure generation strategies.
A comparison with the present approach will close the chapter.
\end{quote}

% ======================================================================
\section{Methodology}\label{sec:methodology}
% ======================================================================

This chapter describes the slicing algorithm as a sequence of
computational stages. Each stage is first motivated by the physical or
robotic requirement it serves, then formalized, and finally discussed
in terms of design decisions and edge cases. The Python realization of
each stage is described in Chapter~\ref{sec:implementation}.

Figure~\ref{fig:pipeline} gives the global picture. The pipeline has two
branches that share the same downstream machinery:

\begin{itemize}
  \item the \textbf{piece branch}: the structural part itself, sliced
        into non-planar layers with fiber-aligned infill;
  \item the \textbf{mold branch}: an automatically generated support
        structure printed before (underneath) the part, sliced into
        near-planar layers with a grid infill.
\end{itemize}

\begin{figure}[htbp]
\centering
\fbox{\parbox{0.92\textwidth}{\centering\vspace{2em}
\textbf{Input:} APDL model (solid mesh + mold surface) + fiber paths (CSV)\\[0.6em]
$\downarrow$\\
\textbf{Stage 1:} model orientation (mean-normal alignment with $+Z$) and lifting\\[0.4em]
$\downarrow$ \hspace{5em} $\downarrow$\\
\textbf{Stage 2:} piece slicing \hspace{2em} \textbf{Stage 2':} shadow volume $\to$ mold slicing\\
(isosurfaces of $q$) \hspace{2.6em} (cone / vertical projection, then $q$-slicing)\\[0.4em]
$\downarrow$ \hspace{5em} $\downarrow$\\
\textbf{Stage 3:} shells + infill per layer (dist-to-boundary field; fiber-aligned planes / fixed XY grid)\\[0.4em]
$\downarrow$\\
\textbf{Stage 4:} path stitching (greedy nearest-endpoint, bounded connections, blockers)\\[0.4em]
$\downarrow$\\
\textbf{Stage 5:} tool-pose computation (normal + tangent $\to$ orthonormal frame)\\[0.4em]
$\downarrow$\\
\textbf{Stage 6:} simplification (RDP / infill collapse) and multi-layer text export\\[0.6em]
\textbf{Output:} trajectory files for RoboDK
\vspace{2em}}}
\caption{Global structure of the slicing pipeline.}
\label{fig:pipeline}
\end{figure}

% ----------------------------------------------------------------------
\subsection{Model preparation: orientation and lifting}\label{sec:meth-orient}
% ----------------------------------------------------------------------

The printing process assumes a dominant build direction: the robot
approaches the part from "above" ($+Z$), and the mold is grown
"downward" from the part's underside onto the printing plane
$Z = z_{\mathrm{plane}}$. A structural model, however, arrives in an
arbitrary orientation in the global frame of the finite-element model.

\subsubsection{Mean-normal estimation and rotation}

The first step therefore estimates the model's natural build direction
and aligns it with $+Z$. The natural build direction is taken to be the
mean of the normals of all layer surfaces of the part. Writing
$\mathcal{S}_k$ for the $k$-th layer surface (their computation is
described in \S\ref{sec:meth-slicing}) and $\mathbf{n}_{k,j}$ for the
unit normal of cell $j$ of $\mathcal{S}_k$, the mean normal is the
surface-area-weighted-free average
\begin{equation}
  \bar{\mathbf{n}} \;=\; \frac{1}{N_{\mathrm{cells}}}
  \sum_{k} \sum_{j} \mathbf{n}_{k,j},
  \qquad N_{\mathrm{cells}} = \sum_k |\mathcal{S}_k|.
\end{equation}
The slicer then builds the rotation matrix $R \in SO(3)$ that maps
$-\bar{\mathbf{n}}$ onto $\mathbf{e}_z = (0,0,1)$:
\begin{equation}
  R = \exp\!\big(\theta\, [\mathbf{a}]_\times\big),
  \qquad
  \mathbf{a} = \frac{\bar{\mathbf{n}} \times \mathbf{e}_z}{\|\bar{\mathbf{n}} \times \mathbf{e}_z\|},
  \quad
  \theta = \arccos\!\big(\bar{\mathbf{n}} \cdot \mathbf{e}_z\big),
\end{equation}
with the degenerate case $\bar{\mathbf{n}} \parallel \mathbf{e}_z$
resolved by $R = I$. The rotation is applied \emph{in place} to all
mesh data of the model: the solid-mesh nodes, the mold mesh, its
Kd-tree, and its PolyData representation, so that every downstream
consumer sees a consistently rotated model.

\subsubsection{Lifting}

After rotation, the model is translated along $+Z$ by an offset so
that its lowest printed surface lies above the printing plane,
leaving room for the support structure. This is a pure translation of
the same data structures, again applied in place.

% ----------------------------------------------------------------------
\subsection{Piece slicing: layers as isosurfaces}\label{sec:meth-slicing}
% ----------------------------------------------------------------------

\subsubsection{Why isosurfaces}

A planar slicer intersects the solid with planes
$z = z_0 + k h$. The generalization adopted here replaces the planes
with the level sets of a scalar field $q$ defined over the volumetric
mesh. The field is chosen such that $q$ is (proportional to) the
signed distance along the local build direction, so that consecutive
integer values of $q$ are one layer height $h$ apart:
\begin{equation}
  \mathcal{S}_k \;=\; \{\,\mathbf{p} \in \mathcal{V} \;:\; q(\mathbf{p}) = q_k\,\},
  \qquad q_{k+1} - q_k = h \ \text{(in physical units)}.
\end{equation}
For a thin-walled part, the isosurfaces of $q$ follow the wall's
mid-surface: the layers are \emph{curved}, and the nozzle deposits
material conformally to the wall. The scalar field itself is produced
upstream by the \texttt{oamc} layer-initialization stage
(\texttt{model.init\_q()}); the slicer only consumes the field named
\texttt{"q"} stored on the volumetric grid, and extracts the level
surfaces with the VTK contour filter.

\subsubsection{Practical extraction}

Each layer surface is extracted with a single contour call on the
unstructured grid, at the level $q_k$. Two properties of the extracted
surface matter for everything downstream:

\begin{itemize}
  \item \textbf{Boundaries.} The isosurface of an open wall is an open
        surface: it has boundary edges (where the contour leaves the
        volume's boundary). These boundary curves are the reference for
        all shell offsets (\S\ref{sec:meth-shells}).
  \item \textbf{Duplicate points.} The contour filter can emit
        duplicate points along the surface boundary, which silently
        corrupts feature-edge detection. The slicer therefore merges
        coincident points (\emph{merge points}, tolerance typically
        0.1\,mm) as the first operation of the path-building stage.
        Because point indexing changes after the merge, every cached
        geometric quantity (normals, distance fields) is invalidated at
        that moment --- the class architecture handles this explicitly
        (see \S\ref{sec:impl-object}).
\end{itemize}

% ----------------------------------------------------------------------
\subsection{The distance-to-boundary field and contour shells}
\label{sec:meth-shells}
% ----------------------------------------------------------------------

\subsubsection{Distance field}

Let $\mathcal{S}$ be one layer surface with point set
$P = \{\mathbf{p}_i\}$ and boundary point set
$B \subset P$ (the points belonging to feature/boundary edges). For
every surface point we define
\begin{equation}
  d(\mathbf{p}) \;=\; \min_{\mathbf{b} \in B} \|\mathbf{p} - \mathbf{b}\|_2,
\end{equation}
the geodesic-free (Euclidean) distance to the layer boundary. The field
is computed exactly with a Kd-tree built on $B$: one parallel
nearest-neighbor query ($O(|P| \log |B|)$) returns all distances at
once. The field is stored as the point-data array
\texttt{dist\_to\_boundary} on the surface, so that all subsequent
VTK scalar-field filters operate on it natively.

\subsubsection{Contour shells}

The shells are the isolines of $d$ inside the layer:
\begin{equation}
  \mathcal{C}_c \;=\; \{\,\mathbf{p} \in \mathcal{S} \;:\; d(\mathbf{p}) = c\,\},
  \qquad c \in \{c_1 < c_2 < \dots\}.
\end{equation}
For the piece, the isovalues are multiples of half a nozzle width
(e.g.\ $c \in \{0.5, 1.0, \dots\}$\,mm with a 0.7\,mm nozzle), giving
perimeter passes at successively larger offsets from the layer edge.
The contour filter produces the isolines; a chaining step then
reconstructs ordered polylines from the filter's segment soup
(\S\ref{sec:meth-chaining}).

\subsubsection{Inner zone and offset surfaces}

The infill must not collide with the shells: it is generated only in
the \emph{inner zone}
\begin{equation}
  \mathcal{Z}_{\mathrm{in}} \;=\; \{\,\mathbf{p} \in \mathcal{S} \;:\; d(\mathbf{p}) \ge c_{\max}\,\},
\end{equation}
where $c_{\max}$ is the largest shell isovalue (the \emph{contour
width}). The zone is obtained by scalar clipping: the VTK clip filter
cuts the surface along the isovalue $c_{\max}$ and re-triangulates the
cut cells, so the boundary of $\mathcal{Z}_{\mathrm{in}}$ lies exactly
on the isoline. (Care with conventions is required: in PyVista's
\texttt{clip\_scalar}, the \texttt{invert} flag selects which side is
kept; the slicer documents explicitly that
\texttt{invert=False} keeps $d \ge c_{\max}$.)

% ----------------------------------------------------------------------
\subsection{Fiber direction and fiber-aligned infill planes}
\label{sec:meth-fibers}
% ----------------------------------------------------------------------

\subsubsection{Mean fiber direction by uncentered PCA}

Each layer carries a set of fiber polylines
$F = \{f_i\}$, each $f_i$ being a sequence of 3D points. The local
stiffness of the printed composite is dominated by the fiber
direction, so the matrix infill must be organized \emph{around} the
fibers. The slicer first estimates the dominant fiber direction of the
layer:

\begin{itemize}
  \item collect the unit direction vectors of every fiber segment,
        skipping zero-length segments:
        $\mathbf{u}_{ij} = (f_i^{j+1} - f_i^j)/\|f_i^{j+1} - f_i^j\|$;
  \item build the uncentered second-moment matrix
        $M = \sum_{i,j} \mathbf{u}_{ij}\mathbf{u}_{ij}^\top$;
  \item the mean direction is the dominant eigenvector of $M$ (largest
        eigenvalue).
\end{itemize}

The absence of centering is deliberate and important: a centered PCA
of a bundle of nearly parallel directions would return the axis of
greatest \emph{scatter}, i.e.\ a direction \emph{perpendicular} to the
fibers. The uncentered moment matrix instead has its dominant
eigenvector aligned with the bundle. The sign of the eigenvector is
fixed deterministically (largest-magnitude component positive) so that
the plane generation is reproducible.

\subsubsection{Infill plane family}

Within the layer surface (of mean unit normal $\mathbf{n}$), the
infill direction is set at $\pm 45^\circ$ from the fiber direction
$\mathbf{d}_f$, alternating with the parity of the layer index:
\begin{equation}
  \mathbf{d}_{\pm} \;=\;
  \cos 45^\circ \, \mathbf{d}_f \;\pm\;
  \sin 45^\circ \, (\mathbf{n} \times \mathbf{d}_f),
\end{equation}
normalized. The alternation gives a cross-ply-like stacking of
successive layers, which balances the bending/torsional stiffness of
the laminate. The infill lines themselves are the intersections of
$\mathcal{Z}_{\mathrm{in}}$ with a family of \emph{parallel planes}
\begin{equation}
  \Pi_s \;=\; \{\,\mathbf{p} : (\mathbf{p} - \mathbf{o}) \cdot \mathbf{m} = s\,\},
  \qquad \mathbf{m} = \mathbf{n} \times \mathbf{d}_\pm,
\end{equation}
spaced by the infill spacing $s_{\mathrm{step}}$ (typically 0.5\,mm
for the piece), centered on the projected extent of the inner zone
with a half-step margin so the extreme planes fully bracket the zone.

% ----------------------------------------------------------------------
\subsection{Mold generation: shadow volumes}\label{sec:meth-shadow}
% ----------------------------------------------------------------------

The part must be printed on a support structure ("mold") that
(a) reaches from the printing plane $Z = z_{\mathrm{plane}}$ up to the
part's lower surface, (b) does not intersect the part, and (c) is
itself printable. The slicer constructs the mold as the
\emph{shadow volume} of the part's lower surface with respect to the
build direction. Three constructions are implemented; all operate on
the part's linearized outer surface (the triangulated skin of the
solid mesh) and all produce a volumetric wedge mesh connecting a "top"
point set (on or near the part) to a "bottom" point set (on the
plane or on the part itself).

\subsubsection{Direct projection}

Every point of the part's downward-facing surface
($\mathbf{n}\cdot\mathbf{e}_z < 0$) is projected vertically onto the
plane. The volume is the set of prisms/wedges between each surface
triangle and its vertical projection. This is the fastest method but
builds material even where the part does not overhang (e.g.\ inside
closed cavities facing up).

\subsubsection{Cone of rays}

For each downward-facing triangle, a cone of rays is cast from its
centroid (offset outward along its normal by a small
\texttt{offset\_faces} to avoid self-intersection): one central
vertical ray and \texttt{no\_rays} rays tilted by the half-angle
\texttt{cone\_angle}, uniformly distributed in azimuth
(\texttt{start\_phase} controlling the rotation of the cone). A
(face, ray) pair is \emph{free} if the ray reaches the printing plane
without intersecting the part's surface. For every free pair, the mold
top point is the (offset) centroid and the mold bottom point is the
ray's intersection with the plane --- so the mold's side walls follow
the cone's slope, giving a self-supporting draft angle. The top points
are triangulated by a 2D Delaunay triangulation of their $XY$
projections, and the wedge volume is built on these triangles.

Faces whose centroid projects too close to the part's silhouette (the
convex hull of the $XY$ projection) are discarded
(\texttt{distance\_from\_contour}), which keeps the mold's rim clean.

\subsubsection{Vertical projection with occlusion}

The most physically accurate variant: each \emph{node} of the
lower surface is projected vertically; if the vertical segment from the
node to the plane crosses the part (which happens for nodes below
overhangs, e.g.\ inside the concavity of a C-beam), the bottom point
is the \emph{first intersection} with the part instead of the plane.
The volume therefore fills the gap between the lower surface and the
nearest surface below it (part or plane). A prism-validity test then
discards the triangles whose extrusion would intersect the part: the
segment between the triangle's top centroid and bottom centroid is
ray-cast against the upward-facing part surface, and any real
intersection inside the segment's height range invalidates the prism.

\subsubsection{Wedge volume assembly}

All three methods share the same volume builder: given top points
$\mathbf{t}_i$, bottom points $\mathbf{b}_i$ and triangles
$(a_j, b_j, c_j)$ of the top surface, the volume is the set of
tetrahedralized wedges connecting each triangle
$(\mathbf{t}_{a}, \mathbf{t}_{b}, \mathbf{t}_{c})$ to its projected
triangle $(\mathbf{b}_{a}, \mathbf{b}_{b}, \mathbf{b}_{c})$. The
result is a \texttt{pv.UnstructuredGrid} of tetrahedra, ready for the
same contour-based slicing as the part.

% ----------------------------------------------------------------------
\subsection{Mold slicing and fixed infill grid}\label{sec:meth-moldslice}
% ----------------------------------------------------------------------

The shadow volume is sliced by the simple planar field
\begin{equation}
  q(\mathbf{p}) = \frac{z(\mathbf{p}) - z_{\mathrm{plane}}}{h},
\end{equation}
so that layer $k$ is the horizontal slice at altitude
$z_{\mathrm{plane}} + k\,h$. For each layer, the pipeline computes:

\begin{enumerate}
  \item \textbf{The layer surface}: the isosurface $q = k$, triangulated
        and with point normals.
  \item \textbf{The distance-to-part field}
        $\mathrm{dist\_to\_part}(\mathbf{p}) = \min\!\big(
        \mathrm{dist}(\mathbf{p}, \text{part surface}),
        \ z(\mathbf{p}) - z_{\mathrm{plane}} \big)$.
        The distance to the (triangulated) part surface is obtained
        with an implicit-distance computation; the minimum with the
        height above the plane closes the field inside cavities, where
        the nearest reference is the build plate rather than a far
        wall of the part.
  \item \textbf{Contour shells}: the isolines of
        $\mathrm{dist\_to\_part}$ at the shell isovalues
        $(i + \tfrac12)\, d_{\mathrm{nozzle}}$, giving two shells
        offset from both the part (above the mold) and the free rim
        of the layer.
  \item \textbf{The inner zone}: the clip
        $\mathrm{dist\_to\_part} \ge n_{\mathrm{shells}}\,
        d_{\mathrm{nozzle}}$.
  \item \textbf{The grid infill}: the intersections of the inner zone
        with the two families of planes $x = x_i$ and $y = y_j$.
\end{enumerate}

\paragraph{Fixed infill levels.}
A subtle but decisive detail: the infill plane positions
$\{x_i\}$, $\{y_j\}$ are computed \textbf{once, from the footprint of
layer zero}, and reused for every layer. If each layer used its own
extent, the planes would shift with the shrinking/growing cross-section
of the mold and the beads would not stack vertically --- the printed
wall would be a lattice of misaligned segments instead of continuous
vertical ribs. With fixed levels, each infill line of layer $k$ lies
exactly on top of the corresponding line of layer $k-1$.

\paragraph{Infill spacing.}
For a two-direction grid in which crossings are printed once, the
area fraction covered by beads of width $w$ spaced by $s$ is
$2w/s - (w/s)^2$ (inclusion--exclusion), hence the spacing for a target
fraction $r$:
\begin{equation}
  s \;=\; \frac{w}{1 - \sqrt{1 - r}}.
\end{equation}
A 100\% infill gives $s = w$ (adjacent beads), 20\% gives
$s \approx 4.9\,w$.

% ----------------------------------------------------------------------
\subsection{Segment chaining}\label{sec:meth-chaining}
% ----------------------------------------------------------------------

Both contour filters (isolines of a scalar field) and plane cuts emit
unordered sets of 2-point line segments. The path logic requires
\emph{ordered polylines}. The slicer chains segments with an explicit
graph traversal rather than relying on VTK's stripper, which was
observed to leave multi-level contour output fragmented:

\begin{itemize}
  \item merge coincident points of the segment soup;
  \item compute the degree of every point; if any point has degree
        $> 2$ (a junction), the graph is not a set of simple chains and
        the chaining is abandoned in favor of a slower, more general
        fallback (one slice call per plane);
  \item otherwise build, for each point, its (at most two) neighbors;
  \item walk every open chain from its degree-1 endpoint, and every
        closed loop from any unvisited degree-2 point (the loop is
        closed explicitly by repeating its first point).
\end{itemize}

The complexity is linear in the number of segments, and the walk
naturally preserves the original sampling of the isolines.

% ----------------------------------------------------------------------
\subsection{Infill crossing management (mold)}\label{sec:meth-crossings}
% ----------------------------------------------------------------------

In a grid infill, the primary and secondary families cross each other.
If both were printed fully, every crossing would be deposited twice
(the second pass would melt/cut through the first bead). The mold
algorithm therefore distinguishes, per layer, a \textbf{primary}
direction (X on even layers, Y on odd layers) from a
\textbf{secondary} direction:

\begin{itemize}
  \item primary lines are printed \emph{continuously};
  \item every secondary line is \textbf{cut at each intersection} with
        a primary line, removing a curvilinear zone of one nozzle
        diameter around the crossing point (half a nozzle width on
        each side), so that the secondary beads approach but never
        overlap the primary beads.
\end{itemize}

The intersections are computed in 2D (XY) between the polyline
segments: for a secondary segment $[\mathbf{p}_1, \mathbf{p}_2]$ and a
primary segment $[\mathbf{q}_1, \mathbf{q}_2]$, the parameters
$(t, u)$ of the intersection solve the $2 \times 2$ cross-product
system
\begin{equation}
  \begin{cases}
    \mathbf{p}_1 + t\,\mathbf{d}_1 \equiv \mathbf{q}_1 + u\,\mathbf{d}_2 \quad (\text{in } XY)\\
    \mathbf{d}_1 = \mathbf{p}_2 - \mathbf{p}_1,\quad
    \mathbf{d}_2 = \mathbf{q}_2 - \mathbf{q}_1
  \end{cases}
\end{equation}
and an intersection exists iff $t, u \in [0,1]$ and the segments are
not parallel. The cut positions are converted to curvilinear
abscissae along the secondary polyline, sorted, merged when
overlapping, and the polyline is split into the surviving
sub-polylines, with exact interpolation of the cut endpoints.

A spatial index (Kd-tree on segment midpoints) restricts the pairwise
tests to segments that are close enough to intersect, turning the
naive $O(NM)$ test into an output-sensitive computation.

\subsubsection{Crossing blocker}

Cutting the secondary lines is not sufficient: the \emph{stitching}
stage (\S\ref{sec:meth-stitching}) may connect the end of one line to
the start of the next with a straight printed move that happens to
cross a primary line --- printing the crossing a second time from
above. The mold path builder therefore attaches a
\textbf{connection blocker} to the stitching algorithm: a callable
that tests the candidate connection segment against all primary
segments with the same 2D intersection test, and forbids the
connection when a crossing is detected. A forbidden connection is
replaced by a leave + approach pair (travel moves), exactly like a
too-long connection.

% ----------------------------------------------------------------------
\subsection{Path stitching}\label{sec:meth-stitching}
% ----------------------------------------------------------------------

Each layer yields a set of curves (shells, infill lines) that must be
printed in a single robotic sequence. The slicer stitches them into
one or several \emph{sub-paths} with the following semantics:

\begin{itemize}
  \item a \textbf{curve} is printed;
  \item a \textbf{connection} is a printed straight move from the end
        of one curve to the start of the next, allowed only if it is
        shorter than a maximum length
        (geometrically derived: the diagonal between two adjacent
        infill lines, $\sqrt{s^2 + h^2}$, times a safety margin) and
        not forbidden by a blocker;
  \item an \textbf{approach} is a travel move (not printed) arriving
        at the start of a sub-path, along the local surface normal,
        at distance \texttt{approach\_distance};
  \item a \textbf{leave} is the symmetric travel move leaving the end
        of a sub-path.
\end{itemize}

\subsubsection{Greedy nearest-endpoint algorithm}

The stitching is a greedy loop over an \emph{endpoint table}: for $N$
curves, a $(2N \times 3)$ array holding the first and last point of
each curve, with an activity mask. Starting from a chosen first curve
(optionally reversed), the loop repeatedly:

\begin{enumerate}
  \item computes the (squared, vectorized) distance from the current
        endpoint to all \emph{active} endpoints;
  \item selects the nearest one, identifies its curve and whether the
        nearest point is the curve's end (in which case the curve is
        consumed in reverse);
  \item checks the connection: length and blocker. If legitimate, a
        connection segment is emitted; otherwise the current sub-path is
        closed (leave move) and a new sub-path is opened (approach
        move) at the chosen curve;
  \item consumes the chosen curve, marks its endpoints inactive, and
        repeats until no curve remains.
\end{enumerate}

The final endpoint gets a leave move. Each approach/leave move is
oriented along the \emph{local surface normal} at the corresponding
curve endpoint, so the tool always arrives at and departs from the
surface perpendicularly --- the same direction it will print in.

\subsubsection{Segment codes}

Every segment of the stitched path carries a type and a command code,
kept distinct between the two material systems of the composite part:

\begin{center}
\begin{tabular}{lcc}
\toprule
Move & Pure (matrix / shells) & Composite (fiber / infill) \\
\midrule
Continuation along a curve & 2 & 0 \\
End of a connection & 3 & 1 \\
Approach (start of sub-path) & 5 & 4 \\
\bottomrule
\end{tabular}
\end{center}

The piece path is built as two fully independent passes --- pure
curves (shells) then composite curves (fiber-aligned infill) --- which
are concatenated in the export. This models the physical process:
the matrix and the fiber tows are deposited by different tool states,
and the robot must finish all matrix material of one family before
switching.

% ----------------------------------------------------------------------
\subsection{Tool-pose computation}\label{sec:meth-frames}
% ----------------------------------------------------------------------

A six-DoF robot needs a full orientation at every trajectory point,
not just a position. The slicer defines the tool frame from two
vector fields along the path:

\begin{itemize}
  \item the \textbf{approach normal} $\mathbf{n}(\mathbf{p})$: the unit
        normal of the surface cell closest to $\mathbf{p}$;
  \item the \textbf{travel tangent} $\mathbf{t}(\mathbf{p})$: the unit
        tangent to the trajectory (central difference at interior
        points, one-sided difference at segment endpoints).
\end{itemize}

The orthonormal frame is then built by Gram--Schmidt:
\begin{equation}
  \mathbf{z} = -\mathbf{n} \ (\text{tool axis pointing into the surface}),
  \quad
  \mathbf{x} = \frac{\mathbf{t} - (\mathbf{t}\cdot\mathbf{z})\,\mathbf{z}}
                     {\|\mathbf{t} - (\mathbf{t}\cdot\mathbf{z})\,\mathbf{z}\|},
  \quad
  \mathbf{y} = \mathbf{z} \times \mathbf{x},
\end{equation}
so that $(\mathbf{x}, \mathbf{y}, \mathbf{z})$ is a right-handed basis,
$\mathbf{x}$ the direction of travel and $\mathbf{z}$ the approach
axis. The degenerate case $\mathbf{t} \parallel \mathbf{z}$ (exactly
the case of approach/leave moves, which travel along the normal) is
handled by an arbitrary perpendicular fallback, since travel
direction carries no meaningful orientation information there.

Two aspects deserve emphasis. First, the tool is \emph{not}
rotationally symmetric about its approach axis: the choice of
$\mathbf{x}$ fixes the tool's rotation about its own axis, and getting
it wrong can mean the tool body colliding with the part even when the
approach axis is correct. Second, the \textbf{mold redefines the
approach normal}: for support material, the natural approach direction
is perpendicular to the \emph{part surface} the support is printed
against, not to the support layer's own surface. The mold class
therefore overrides the normal lookup to sample the part's surface
normals instead --- one of the two places where the class hierarchy
changes behavior rather than merely adding it.

Tangents are computed independently within each segment, so a small
discontinuity in the frame's $\mathbf{x}$ axis can occur exactly at a
segment boundary (curve $\to$ connection): only the rotation of the
tool about its approach axis is affected, never the position, and the
robot interpolates smoothly over it.

% ----------------------------------------------------------------------
\subsection{Polyline simplification}\label{sec:meth-simplify}
% ----------------------------------------------------------------------

Raw isolines can contain thousands of nearly collinear points. The
export stage offers two levels of simplification:

\begin{itemize}
  \item \textbf{Ramer--Douglas--Peucker (RDP)} \cite{rdp}, implemented
        iteratively (explicit stack, no recursion) and directly in 3D:
        the perpendicular distance of each candidate point to the
        current 3D chord is computed exactly; points are kept as long
        as the maximum deviation exceeds the tolerance. Endpoints are
        always kept, and simplification is applied per segment, so
        distinct segments are never merged.
  \item \textbf{Infill collapse} (mold): a mold infill line is, in the
        normal case, the intersection of a planar layer with a plane
        $x = \text{const}$ or $y = \text{const}$ --- a straight line.
        It can therefore be reduced to its two endpoints, shrinking
        the trajectory file by orders of magnitude. A safety net
        guards this collapse: the line is flattened only if every
        point lies within a tolerance of the endpoint chord \emph{and}
        the points progress monotonically along the chord. A line that
        doubles back (contains a printed jump) is left untouched, and
        a merely curved line falls back to RDP; no gap is ever erased.
\end{itemize}

% ----------------------------------------------------------------------
\subsection{Trajectory export}\label{sec:meth-export}
% ----------------------------------------------------------------------

The final trajectory is a plain-text table, one row per trajectory
point:
\begin{center}
\texttt{x y z ux uy uz vx vy vz wx wy wz type [layer]}
\end{center}
where $(u\!\cdot\!\cdot, v\!\cdot\!\cdot, w\!\cdot\!\cdot)$ are the
three columns of the rotation matrix $R$ (the tool's X, Y and Z axes
in world coordinates), \texttt{type} is the segment label
(\texttt{approach}, \texttt{curve}, \texttt{connection},
\texttt{leave}) and \texttt{layer} the layer index. A small
post-processing drops the duplicated point at each segment junction
(consecutive segments share one point). Multi-layer export writes the
header once and appends every subsequent layer to the same file, in
print order.

This format is deliberately minimal and self-describing: building the
$4\times 4$ pose matrix on the receiving end is a reshape, and the
type labels let the robot-side post-processor map moves to printed /
travel motions and select the deposition tool state.

% ======================================================================
\section{Implementation}\label{sec:implementation}
% ======================================================================

This chapter describes the Python realization of the methodology of
Chapter~\ref{sec:methodology}: the package layout, the class
architecture, the role of the main methods with their signatures and
the non-trivial implementation choices, the performance engineering,
and the connection to the robot simulator RoboDK.

The package is small by design --- about 2\,750 lines of Python spread
over three modules --- because all heavy geometry is delegated to VTK
filters through PyVista, and all array mathematics is vectorized with
NumPy. The implementation effort concentrates on three things:
\emph{correctness} of the geometric conventions (clipping, normals,
chaining), \emph{performance} of the hot paths (caching and
vectorization), and \emph{extensibility} (the class hierarchy
separating generic layer behavior from piece- and mold-specific
behavior).

\begin{table}[htbp]
\centering
\begin{tabular}{lll}
\toprule
Module & Lines & Content \\
\midrule
\texttt{src/slicer/SlicedObject.py} & 687 & base class: geometry, path stitching, frames, export \\
\texttt{src/slicer/SlicedPiece.py}  & 475 & composite piece layers: fibers, slicing planes, model prep \\
\texttt{src/slicer/SlicedMold.py}   & 1589 & mold: shadow volumes, mold slicing, infill crossing, paths \\
\bottomrule
\end{tabular}
\caption{Modules of the \texttt{slicer} package.}
\label{tab:modules}
\end{table}

% ----------------------------------------------------------------------
\subsection{Package layout and dependencies}
% ----------------------------------------------------------------------

The package follows the \texttt{src/} layout with a minimal
\texttt{pyproject.toml} (setuptools backend, Python $\ge$ 3.10):
\begin{lstlisting}[caption={Package layout.}]
slicer/
  pyproject.toml
  src/slicer/
    __init__.py        # exports SlicedObject, SlicedPiece, SlicedMold
    SlicedObject.py    # base class
    SlicedPiece.py     # the composite part
    SlicedMold.py      # the support structure
  examples/
    print_layers.py    # end-to-end example: piece
    mold_generation.py # end-to-end example: mold + fibers
    fibers/layer_*/    # fiber path CSV files
    c_beam_large_thin.dat  # APDL model
\end{lstlisting}
Dependencies are \texttt{pyvista}, \texttt{numpy} and \texttt{scipy};
the slicer additionally imports the model layer from the companion
library \texttt{oamc} (\texttt{CompositeModel},
\texttt{compute\_level\_surface}) and, on the mold side, a few VTK
classes directly (\texttt{vtkStaticCellLocator}, \texttt{vtkPoints},
\texttt{vtkIdList}) for the performance-critical ray casting
(\S\ref{sec:impl-perf}).

% ----------------------------------------------------------------------
\subsection{Class architecture}\label{sec:impl-arch}
% ----------------------------------------------------------------------

\begin{figure}[htbp]
\centering
\fbox{\parbox{0.9\textwidth}{\centering\vspace{1em}
\texttt{SlicedObject}\\
\footnotesize surface + caches; dist-to-boundary; offset zones;\\
\footnotesize contour polylines; local normals; segment chaining;\\
\footnotesize single-pass stitching; tool frames; RDP; text export\\[0.8em]
$\uparrow$\hspace{0.6em}$\uparrow$\\[0.2em]
\texttt{SlicedPiece}\hspace{2em}\texttt{SlicedMold}\\
\footnotesize fiber paths, mean direction,\hspace{2.2em}\footnotesize shadow volumes, dist-to-part, fixed grid\\
\footnotesize fiber-aligned planes, model rotation\hspace{0.4em}\footnotesize infill crossing cuts, crossing blocker,\\
\footnotesize \& lifting, isosurface layers\hspace{3.7em}\footnotesize overridden approach normals, mold paths
\vspace{1em}}}
\caption{Class hierarchy of the \texttt{slicer} package.}
\label{fig:classes}
\end{figure}

The three classes split responsibilities along one axis: \emph{what is
generic to any sliced layer} versus \emph{what is specific to what is
being printed}.

\begin{itemize}
  \item \texttt{SlicedObject} is instantiated with a single
        \texttt{pv.PolyData} --- one layer surface --- and knows nothing
        about the object the surface comes from. It provides the
        distance-to-boundary machinery, offset zones, contour
        extraction, polyline handling, the entire path-stitching
        engine, tool-frame computation, simplification and export.
  \item \texttt{SlicedPiece} \emph{adds} fiber knowledge: it carries a
        dictionary of fiber polylines per layer, computes their mean
        direction, generates the fiber-aligned infill planes, and adds
        the model-level preparation (rotation matrix, layer extraction
        from the \texttt{oamc} grid).
  \item \texttt{SlicedMold} both \emph{adds} (shadow-volume
        construction, distance-to-part field, fixed XY grid, crossing
        management) and \emph{overrides}: its approach/leave moves use
        the \emph{part} surface normals instead of the mold layer
        normals.
\end{itemize}

The override deserves a comment, because it is the hinge of the
design. The base class exposes two normal lookups,
\texttt{get\_local\_normal} (single point) and
\texttt{get\_local\_normals} (batched). All downstream consumers ---
approach/leave construction, frame computation, path export --- go
through the batched lookup, and the batched lookup of
\texttt{SlicedMold} samples the part surface. To make single-point and
batched behavior impossible to diverge, the base class routes
\texttt{get\_local\_normal} through \texttt{get\_local\_normals} and
provides an internal dispatcher (\texttt{\_normals\_at}) that calls the
per-point method \emph{only if} a subclass has overridden it, and the
fast batched method otherwise. A subclass therefore only needs to
override the batched method to change the normal semantics everywhere
at once.

% ----------------------------------------------------------------------
\subsection{\texttt{SlicedObject}: the generic layer}\label{sec:impl-object}
% ----------------------------------------------------------------------

\subsubsection{Surface property and cache invalidation}

The class holds its geometry in a \texttt{surface} property with a
setter. The setter is not cosmetic: assigning a new mesh invalidates
two cached quantities, the cell-normals array and the
distance-to-boundary state flag. This matters because the piece path
builder \emph{reassigns} \texttt{self.surface} mid-algorithm (after
merging duplicate points, which renumbers all points); without the
invalidation, stale normals and distance fields would silently poison
every subsequent computation. The caching is lazy: the cell normals,
for instance, are normalized once on first access and stored; PyVista's
\texttt{cell\_normals} attribute recomputes the whole mesh when no
normals array is attached, so an accidental per-point access in a loop
would otherwise dominate the runtime.

\subsubsection{Distance-to-boundary and offset zones}

\begin{lstlisting}[caption={Distance field and offset zone (essence).}]
def build_distance_tree(self):
    edges = self.surface.extract_feature_edges()
    boundary_points = edges.points
    tree = cKDTree(boundary_points)
    distances, _ = tree.query(self.surface.points, workers=-1)
    self.surface["dist_to_boundary"] = distances

def build_offset_surface(self, value, invert=False):
    if not self._distance_tree_built:
        self.build_distance_tree()
    return self.surface.clip_scalar(scalars="dist_to_boundary",
                                    value=value, invert=invert)
\end{lstlisting}
Boundary detection is VTK's feature-edge extraction on the layer
surface; the nearest-neighbor query runs on all cores
(\texttt{workers=-1}). The offset zone is a scalar clip: VTK cuts and
re-triangulates every cell crossed by the isovalue, so the zone's
boundary lies \emph{exactly} on the isoline $d = c$. The docstring
spells out the \texttt{invert} convention (\texttt{invert=True} keeps
$d \le c$), which is the single most error-prone convention of the
whole pipeline.

\subsubsection{Contour polylines and segment chaining}

Isolines are extracted with the contour filter on
\texttt{dist\_to\_boundary} and converted to ordered polylines. The
conversion goes through \texttt{vtkStripper} (via
\texttt{PolyData.strip}) and a manual walk over the stripped
connectivity array: VTK stores polylines as
\texttt{[n$_1$, id$_{11}$\,\dots\,id$_{1n_1}$, n$_2$, \dots]}, and the
parser collects each run as a point array. A second, more robust
chaining (\texttt{\_chain\_segments}, used for plane cuts) works
directly on the raw segment graph: degree counting, junction detection,
neighbor walks --- described in \S\ref{sec:meth-chaining}.

\subsubsection{Single-pass stitching}

The core of the path logic is \texttt{build\_path\_single\_pass},
which implements the greedy algorithm of \S\ref{sec:meth-stitching}.
Its implementation is dominated by two data structures:

\begin{itemize}
  \item the \emph{endpoint table}: a $(2N \times 3)$ NumPy array with
        row $2k$ the first point of curve $k$ and row $2k+1$ its last
        point, plus a boolean activity mask. The nearest-endpoint query
        is then a single vectorized squared-distance computation
        (\texttt{einsum}) with masked-out inactive rows --- the
        previous implementation rebuilt a Kd-tree at every iteration,
        which is exactly the kind of per-iteration setup cost the
        rewrite removes.
  \item the \emph{command list} \texttt{Lcmd}: for each printed
        endpoint, the coordinate and the code
        (approach/connection/continuation), later consumed by the
        export.
\end{itemize}

Sub-paths are delimited by \texttt{approach}/\texttt{leave} segments
built along the endpoint normals (batched, from the endpoint table).
The optional \texttt{connection\_blocker} --- a callable
$(p_{\text{start}}, p_{\text{end}}) \to \texttt{bool}$ --- lets the
caller inject a topological constraint into the stitching without the
stitching engine knowing anything about it; the mold passes its
crossing blocker there (\S\ref{sec:impl-mold}).

\texttt{build\_full\_printing\_path} runs the engine twice --- once on
the pure curves (shells) with codes (5, 3, 2), once on the composite
curves (fiber-aligned infill) with codes (4, 1, 0) --- and concatenates
the two passes.

\subsubsection{Tool frames and export}

The frame machinery mirrors the methodology of
\S\ref{sec:meth-frames}:
\begin{itemize}
  \item \texttt{compute\_path\_normals} stacks all points of all
        segments into one array and issues a \emph{single} batched
        closest-cell query (the locator is built once internally);
  \item \texttt{compute\_path\_tangents} uses central differences
        inside each segment and one-sided differences at the ends;
  \item \texttt{build\_orthonormal\_frames} is fully vectorized over
        all points of all segments: Gram--Schmidt against the
        (sign-flipped) normals, degenerate-tangent fallback, and a
        final \texttt{np.split} back into per-segment arrays. The
        output convention is documented in the docstring: column 0/1/2
        of each $3\times3$ frame is the tool X/Y/Z axis in world
        coordinates, i.e.\ the rotation part of a robot pose.
\end{itemize}

The export path is \texttt{export\_printing\_path}, a convenience
wrapper that (optionally) simplifies each segment with RDP
(\texttt{simplify\_polyline}, the iterative implementation described in
\S\ref{sec:meth-simplify}), drops the duplicated junction points
(\texttt{drop\_shared\_points}), computes the orientations \emph{on the
simplified points} --- so the written orientation always matches the
written geometry --- and delegates to \texttt{export\_path\_to\_txt}.
The low-level writer builds one printf-style format string per segment
(baking in the constant type/layer labels) and writes each segment with
a single \texttt{np.savetxt} call; in \texttt{"a"} mode it appends
rows without rewriting the header, which is how multi-layer export
accumulates one file.

% ----------------------------------------------------------------------
\subsection{\texttt{SlicedPiece}: the composite part}\label{sec:impl-piece}
% ----------------------------------------------------------------------

\subsubsection{Fiber data}

\texttt{load\_fiber\_paths} reads the CSV files of a layer directory
(\texttt{fibers/layer\_\textit{L}/fiber\_\textit{L}\_\textit{i}.csv}),
each file one polyline. The mean direction is computed once per
instance and cached; when \texttt{slice\_model} builds all layers with
the same fiber dictionary, the cache of the first layer is copied to
the others, so the PCA runs exactly once per print job.

The PCA itself is three lines of NumPy --- directions stacked, uncentered
second-moment matrix, dominant eigenvector --- but it is guarded by two
documented subtleties: the centering trap (\S\ref{sec:meth-fibers}) and
the sign determinism.

\subsubsection{Slicing planes and single-pass cutting}

\texttt{generate\_slicing\_planes} implements the $\pm45^\circ$
direction of \S\ref{sec:meth-fibers} and returns the plane family as a
list of \texttt{\{"origin", "normal"\}} dictionaries.

The cutting is where a major performance idea lives:
\texttt{\_cut\_with\_planes} does \emph{not} call \texttt{slice} once
per plane. Since the planes are parallel, they are the level sets of the
single scalar field $s(\mathbf{p}) = \mathbf{p}\cdot\mathbf{m}$; the
method attaches that field to the zone and calls the contour filter
once with \emph{all} the levels. One mesh traversal replaces $n$. The
chaining then runs on the multi-level contour output, with the
junction-fallback of \S\ref{sec:meth-chaining} (the fallback is
\texttt{\_cut\_with\_planes\_slow}, one \texttt{slice} per plane, used
only when the graph has junctions).

\subsubsection{Model-level preparation}

\begin{itemize}
  \item \texttt{\_layer\_levels} queries the \texttt{oamc} grid for the
        list of $q$ values; \texttt{get\_printing\_layers} extracts the
        isosurfaces (optionally only a subset --- extracting isosurfaces
        is the expensive model-level operation); \texttt{slice\_model}
        returns a list of \texttt{SlicedPiece};
        \texttt{from\_model} builds a single layer by index, extracting
        only the requested isosurface.
  \item \texttt{get\_mean\_surface\_normals} averages cell normals over
        all layer surfaces; \texttt{get\_rotation\_matrix} builds the
        axis--angle rotation of \S\ref{sec:meth-orient} with
        \texttt{scipy.spatial.transform.Rotation}.
\end{itemize}

\subsubsection{Per-layer path assembly}

\texttt{build\_printing\_path} chains the operations of
\S\ref{sec:meth-shells}--\ref{sec:meth-fibers} on one layer: merge
points (with the cache invalidation this triggers), generate the
planes, clip the inner zone, cut the infill in one pass, extract the
contour-shell isolines, and stitch the two passes. The parameters
expose the physically meaningful knobs: contour width (3\,mm in the
example), approach distance (1\,mm), max connection length, and the
shell isovalues.

\texttt{export\_layers\_to\_txt} loops over all layers, builds each
path and appends to the trajectory file
(\texttt{mode="w"} for the first layer, \texttt{"a"} afterwards).

% ----------------------------------------------------------------------
\subsection{\texttt{SlicedMold}: the support structure}\label{sec:impl-mold}
% ----------------------------------------------------------------------

The mold module is the largest of the three; it contains the shadow
volume builders, the mold slicing loop, the infill crossing machinery,
and the mold path builder.

\subsubsection{Overridden normals and the \texttt{\_RayCaster} helper}

The class stores the \emph{part} surface (its own property, with the
same cache-invalidation pattern as the base class) and overrides
\texttt{get\_local\_normals} to query the part's closest cells ---
implementing the "approach against the part" semantics of
\S\ref{sec:meth-frames}.

The module also defines \texttt{\_RayCaster}, a thin wrapper around
\texttt{vtkStaticCellLocator}: the locator is built \emph{once} per
surface, and \texttt{hits(start, end)} returns every intersection of a
segment with the mesh. The wrapper exists because
\texttt{PolyData.ray\_trace} may rebuild its locator on every call
depending on the PyVista version --- and the shadow-volume builders
cast thousands of rays.

\subsubsection{Three shadow-volume builders}

\begin{itemize}
  \item \texttt{get\_shadow\_volume}: extract the downward-facing cells
        of the outer surface, project their points vertically, build
        the wedges (\texttt{\_wedge\_volume}).
  \item \texttt{get\_shadow\_volume\_cone}: the cone-of-rays method of
        \S\ref{sec:meth-shadow}. The implementation is a textbook case
        of the project's performance doctrine: for every ray direction,
        all \emph{cheap} tests (does the ray even point down? is the
        plane hit inside the part's bounding box? is it far enough from
        the convex-hull contour?) are evaluated for all faces at once
        as boolean NumPy masks; only the surviving (face, ray) pairs
        are actually cast against the single locator. The free pairs
        are then assembled by a 2D Delaunay triangulation of the top
        points and the wedge builder.
  \item \texttt{get\_shadow\_volume\_vertical}: the occlusion-aware
        vertical projection. Per node, the bottom point defaults to the
        vertical projection onto the plane and is replaced by the
        first intersection with the part's upward-facing surface when
        the vertical segment is blocked; per triangle, a prism-validity
        ray test discards wedges whose extrusion would cross the part.
        Again, bounding-box pre-filters keep the actual ray count low,
        and the method also returns the two point clouds for
        visualization.
\end{itemize}

\texttt{slice\_model} composes the vertical projection with the
slicing loop in one call; \texttt{rotate\_model} /
\texttt{lift\_model} implement the in-place model transformation of
\S\ref{sec:meth-orient} (nodes, Kd-tree and PolyData of the mold are
all updated, using \texttt{object.\_\_setattr\_\_} to bypass the frozen
dataclass fields of the \texttt{oamc} model).

\subsubsection{The slicing loop}

\texttt{slice\_shadow\_volume} implements \S\ref{sec:meth-moldslice}.
Beyond what the methodology chapter describes, three implementation
details are worth recording:

\begin{itemize}
  \item the \textbf{fixed infill levels} are computed from the
        \emph{bounds of layer zero's isosurface} (falling back to the
        volume bounds if the first layer is empty) and stored once on
        the class, then passed to every \texttt{SlicedMold} instance
        --- the fix that earned the project's first pull request;
  \item the \textbf{distance-to-part} uses
        \texttt{compute\_implicit\_distance} against the triangulated,
        cleaned part surface, rather than sampling the volume's field
        (the contour filter's carried values differ from the volume's
        own interpolation);
  \item the shells are chained through \texttt{vtkStripper} here
        (single-level contours, where the stripper behaves), unlike
        the multi-level piece contours where manual chaining is
        required.
\end{itemize}

The grid cutting itself reuses the single-pass idea in its axis form
(\texttt{\_cut\_axis\_planes}): the coordinate
$x$ (resp.\ $y$) is the scalar field, all levels are contoured at once,
chained, and the polylines are then re-sorted plane by plane.

\subsubsection{Infill crossing machinery}

The crossing management of \S\ref{sec:meth-crossings} is factored
around a shared engine, \texttt{\_cut\_secondary\_lines}, with two
public wrappers:

\begin{itemize}
  \item \texttt{cut\_infill\_at\_intersections}: exact
        version, every (secondary, primary) segment pair tested;
  \item \texttt{cut\_infill\_at\_intersections\_fast}: a Kd-tree on
        primary-segment midpoints restricts each secondary segment to
        the primary segments within
        (half its length + spacing + nozzle width) --- the tests
        themselves stay vectorized per secondary polyline.
\end{itemize}

The curvilinear bookkeeping is generic polyline code:
\texttt{interpolate\_polyline\_at} (binary search in the cumulative
length, linear interpolation), \texttt{sample\_polyline\_between}, and
\texttt{\_cut\_polyline\_at\_positions} (exclusion zones of radius
$d_{\mathrm{nozzle}}/2$ around each cut, merged when overlapping).

\subsubsection{Mold path building}

\texttt{build\_support\_path} assembles one layer:

\begin{enumerate}
  \item computes the geometric max connection length
        (\texttt{compute\_max\_connection\_length}: the diagonal
        $\sqrt{s^2 + h^2}$ with a 1.1 safety margin);
  \item sorts each infill family into \emph{zigzag order}
        (by mean coordinate along the family axis, every other line
        reversed) --- this is what makes the greedy stitching produce
        the classic back-and-forth sweep instead of random jumps;
  \item assigns primary/secondary by layer parity and cuts the
        secondary family at the crossings;
  \item optionally simplifies: RDP on the shells
        (\texttt{shell\_simplify\_tolerance}) and/or infill collapse
        (\texttt{collapse\_infill}, guarded by
        \texttt{collapse\_infill\_line} and its monotonicity safety
        net, \S\ref{sec:meth-simplify});
  \item stitches two passes with \texttt{build\_path\_single\_pass}:
        shells first, then infill (primary + cut secondary) with the
        \textbf{crossing blocker}
        (\texttt{\_make\_crossing\_blocker}: a closure that broadcasts
        the candidate connection against all primary segments through
        the vectorized 2D intersection test --- a forbidden crossing
        becomes a leave + approach pair).
\end{enumerate}

\texttt{export\_layers\_to\_txt} loops over all mold layers and writes
the single trajectory file, exactly like the piece version.

% ----------------------------------------------------------------------
\subsection{Performance engineering}\label{sec:impl-perf}
% ----------------------------------------------------------------------

Performance was a first-class requirement, not an afterthought: the
input meshes are finite-element models with tens of thousands of cells,
and the naive implementation of several stages is quadratically slow.
The doctrine applied throughout the package is uniform and worth
summarizing:

\begin{enumerate}
  \item \textbf{Compute once, cache.} Every quantity that depends only
        on the mesh (cell normals, Kd-trees, locators, PCA of the
        fibers) is computed on first use and stored; caches are
        explicitly invalidated when the geometry changes (the
        \texttt{surface} setter) or shared between instances (the
        fiber direction across layers).
  \item \textbf{Batch every query.} All closest-cell lookups, distance
        queries and normal queries are issued for entire point arrays
        in one call, never point by point; the locator build cost is
        amortized over the whole batch. The stitching's
        nearest-endpoint search is a masked vectorized distance
        computation, replacing a per-iteration Kd-tree build.
  \item \textbf{Replace repeated traversals by single multi-level
        filters.} Parallel plane cuts are one contour call with all
        levels (\S\ref{sec:impl-piece}); axis-plane grids are one
        contour call per direction (\S\ref{sec:impl-mold}).
  \item \textbf{Filter before you cast.} In the ray-based stages, every
        test that does not need an actual ray cast (bounding boxes,
        distance to the convex hull, direction signs) is evaluated for
        all faces at once as NumPy masks; only survivors reach the
        single locator. One of the builders' docstrings states it as a
        rule: "only the surviving (face, ray) pairs are actually
        cast".
  \item \textbf{Index pairwise tests.} The infill crossing tests use a
        Kd-tree on segment midpoints to avoid $O(NM)$ pairwise
        segment intersection.
  \item \textbf{Vectorize I/O.} The trajectory writer builds per-segment
        format strings and writes each segment with one
        \texttt{np.savetxt} call, instead of formatting every point in
        Python.
\end{enumerate}

% ----------------------------------------------------------------------
\subsection{Interface with RoboDK}\label{sec:impl-robodk}
% ----------------------------------------------------------------------

The slicer never talks to the robot directly; the contract between
them is the exported trajectory file, and RoboDK is the offline
programming environment where the trajectory becomes a robot program.

\subsubsection{What the export provides}

For every trajectory point, RoboDK receives:
\begin{itemize}
  \item the position $(x, y, z)$;
  \item the full orientation, as the nine components of the tool frame
        --- directly reshaped into the rotation block of a $4 \times 4$
        pose matrix;
  \item the segment type, which distinguishes \emph{printed} moves
        (\texttt{curve}, \texttt{connection}) from \emph{travel} moves
        (\texttt{approach}, \texttt{leave}) --- the equivalent of
        G-code's distinction between extruding and travel moves, at
        robot level;
  \item the pass family (pure vs.\ composite command codes, 
        \S\ref{sec:meth-stitching}), telling the post-processor which
        deposition state to select;
  \item the layer index, so the program can be organized into layers
        and resumed per layer.
\end{itemize}

This is strictly more information than a planar slicer provides, and
each item exists for a robotic reason: the orientation because a
six-DoF arm has three rotational DoF to feed; the travel/printed
distinction because approach and leave moves run at travel speed
without extrusion; the layer index because the physical process
prints layer by layer even when the layers are curved.

\subsubsection{RoboDK-side workflow}

The intended RoboDK workflow is:
\begin{enumerate}
  \item import the trajectory file (a Python or CSV post-processor
        reads each row, builds the pose from the coordinates and
        rotation block, and emits a \texttt{MoveJ} to the first
        approach point and \texttt{MoveL} targets along the path);
  \item attach the deposition tool (extruder) to the robot and set the
        tool frame so that the trajectory's approach axis
        ($\mathbf{z}$, pointing into the surface) matches the tool's
        spindle axis;
  \item simulate: RoboDK checks reachability and collisions against
        the cell and the part being built, which the slicer
        deliberately does not do (\S\ref{sec:intro});
  \item generate the program for the target robot controller.
\end{enumerate}

The \texttt{robodk} Python package is listed in the project
dependencies (commented out in \texttt{pyproject.toml}, since it is
only needed in the robot cell, not for slicing), and the exported
row format (position + 9 rotation components + labels) maps one-to-one
onto RoboDK's \texttt{AddTarget}/pose APIs.

\subsubsection{Why not G-code}

A classical slicer would export G-code; this project exports a pose
table instead, for three reasons. First, G-code's rotary state (A/B/C
or I/J/K angles) is machine-convention-dependent, while a rotation
matrix is unambiguous --- the slicer stays agnostic to the robot brand.
Second, the path semantics here (approach along the surface normal,
two material families, non-planar layers) do not fit the planar
G-code model of "layers at fixed Z". Third, RoboDK's offline
programming model is exactly designed to consume generic pose
targets and translate them per controller, which removes the need for
the slicer to know anything about the robot's kinematics.

% ----------------------------------------------------------------------
\subsection{Example programs}\label{sec:impl-examples}
% ----------------------------------------------------------------------

The repository ships two executable end-to-end examples
(\texttt{examples/}), which double as integration tests of the whole
pipeline:

\begin{itemize}
  \item \texttt{print\_layers.py}: parses the C-beam APDL model, builds
        the \texttt{CompositeModel} with its matrix and fiber
        materials, rotates and lifts the model, loads the fiber CSVs,
        slices the model into \texttt{SlicedPiece} layers, visualizes
        two layers (with an exaggerated visual layer spacing for
        readability --- the real spacing is 0.35\,mm) with the path
        segments color-coded by type, and exports the piece trajectory
        with orientations to
        \texttt{trajectories/piece\_path.txt}.
  \item \texttt{mold\_generation.py}: same model preparation, then the
        \texttt{oamc} structural pipeline (compliance evaluation,
        fiber field computation, RDP downsampling, fiber export to
        CSV) and the mold branch: shadow-volume generation and
        slicing into \texttt{SlicedMold} layers, with the mold
        trajectory export.
\end{itemize}

Together they demonstrate the intended division of labor:
\texttt{oamc} decides \emph{where the fibers must go} (structural
optimization), and \texttt{slicer} decides \emph{how to print them}
(geometry and trajectories).

% ======================================================================
\section{Results}\label{sec:results}
% ======================================================================

% ----------------------------------------------------------------------
% NOTE: this section is intentionally left as a placeholder.
% It will be completed later (print quality, trajectory statistics,
% RoboDK simulation screenshots, etc.).
% ----------------------------------------------------------------------

\begin{quote}
\itshape
\textbf{[To be completed.]} This chapter will present the results of
the pipeline on the C-beam case study: the generated non-planar layers
and fiber-aligned infill, the shadow volume and mold layers, statistics
of the generated trajectories (number of segments, travel distance,
simplification gains), and RoboDK simulation screenshots of the
six-DoF toolpaths.
\end{quote}


\bibliographystyle{plain}
\begin{thebibliography}{9}
  \bibitem{pyvista} The PyVista Developers, \emph{PyVista: 3D plotting and mesh analysis through a streamlined, VTK-based interface}, \url{https://docs.pyvista.org}.
  \bibitem{vtk} Kitware, \emph{The Visualization Toolkit (VTK)}, \url{https://vtk.org}.
  \bibitem{scipy} P.~Virtanen et al., \emph{SciPy 1.0: Fundamental Algorithms for Scientific Computing in Python}, Nature Methods 17 (2020).
  \bibitem{rdp} D.~Douglas and T.~Peucker, \emph{Algorithms for the reduction of the number of points required to represent a digitized line or its caricature}, Cartographica 10(2), 1973.
  \bibitem{robodk} RoboDK, \emph{RoboDK documentation}, \url{https://robodk.com/doc}.
\end{thebibliography}

\end{document}
